\clearpage
\section*{常用 \LaTeX{} 示例代码（赛前学习，提交前删除）}
本节只用于赛前熟悉写法，正式提交前请整体删除。本节中的图片示例写在代码框里，不会真的读取图片文件；正式使用时，把图片放入 \texttt{figures/} 文件夹，再把对应代码复制到正文位置。

\subsection*{1. 大标题、小标题}
\begin{tcode}
\title{论文题目}

\section{一级标题}
这里写正文。

\subsection{二级标题}
这里写更细的说明。

\subsubsection{三级标题}
三级标题不要过多使用，除非结构确实需要。
\end{tcode}

\subsection*{2. 单张图片}
\begin{tcode}
\begin{figure}[!htbp]
  \centering
  \includegraphics[width=0.75\textwidth]{figures/result1.png}
  \caption{图片标题}
  \label{fig:result1}
\end{figure}

如图~\cref{fig:result1} 所示，模型结果呈现出明显变化。
\end{tcode}

\subsection*{3. 并排图片}
\begin{tcode}
\begin{figure}[!htbp]
  \centering
  \begin{minipage}[c]{0.48\textwidth}
    \centering
    \includegraphics[width=\textwidth]{figures/result_a.png}
    \subcaption{方案 A}
    \label{fig:result-a}
  \end{minipage}
  \hfill
  \begin{minipage}[c]{0.48\textwidth}
    \centering
    \includegraphics[width=\textwidth]{figures/result_b.png}
    \subcaption{方案 B}
    \label{fig:result-b}
  \end{minipage}
  \caption{不同方案的结果对比}
  \label{fig:compare}
\end{figure}

图~\cref{fig:compare} 给出了整体对比，其中 \cref{fig:result-a} 为方案 A。
\end{tcode}

\subsection*{4. 公式与引用}
\begin{tcode}
行内公式写法：目标函数 $f(x)$ 用于衡量方案优劣。

不编号公式：
\[
  z = \sum_{i=1}^{n} c_i x_i
\]

带编号公式：
\begin{equation}
  \min Z = \sum_{i=1}^{n} c_i x_i
  \label{eq:objective}
\end{equation}

由式~\cref{eq:objective} 可得到最优目标值。

多行对齐公式：
\begin{align}
  s_i &= x_i - \bar{x},\\
  \sigma &= \sqrt{\frac{1}{n}\sum_{i=1}^{n}s_i^2}.
\end{align}
\end{tcode}

\subsection*{5. 三线表}
\begin{tcode}
\begin{table}[!htbp]
  \centering
  \caption{不同模型的评价结果}
  \label{tab:model-compare}
  \begin{tabularx}{\textwidth}{cXXX}
    \toprule
    模型 & 准确率 & 运行时间 & 说明 \\
    \midrule
    模型一 & 0.912 & 12.5 s & 基准模型 \\
    模型二 & 0.936 & 18.7 s & 综合表现较好 \\
    模型三 & 0.928 & 10.1 s & 速度较快 \\
    \bottomrule
  \end{tabularx}
\end{table}

从表~\cref{tab:model-compare} 可以看出，模型二的准确率最高。
\end{tcode}

\subsection*{6. 参考文献引用}
\begin{tcode}
正文中引用文献：灰色预测模型常用于小样本时间序列预测~\cite{ref:grey}。

\begin{thebibliography}{99}
  \bibitem{ref:grey}
  邓聚龙. 灰色系统理论教程[M]. 武汉: 华中理工大学出版社, 1990.

  \bibitem{ref:data}
  国家统计局. 中国统计年鉴[M]. 北京: 中国统计出版社, 2025.
\end{thebibliography}
\end{tcode}

\subsection*{7. 代码块}
\begin{tcode}
\begin{lstlisting}[language=Python]
import numpy as np

def normalize(x):
    x = np.asarray(x, dtype=float)
    return (x - x.min()) / (x.max() - x.min())

scores = normalize([82, 91, 76, 88])
print(scores)
\end{lstlisting}
\end{tcode}
